{\fontsize{9}{11}\selectfont
\begin{longtable}{p{4.6cm}p{2.6cm}p{6.1cm}}
  \caption{Parametros dinamicos del muro exterior base} \\
  \toprule
  Indicador & Valor & Interpretacion \\
  \midrule
  \endfirsthead
  \caption[]{Parametros dinamicos del muro exterior base (continuacion)} \\
  \toprule
  Indicador & Valor & Interpretacion \\
  \midrule
  \endhead
  \bottomrule
  \endfoot
  Componente evaluado & Enlucido de yeso 15 mm + bloque de hormigon 150 mm + enlucido de cemento 15 mm & Corresponde al muro exterior base equivalente a la construccion \texttt{01.WA.ENL1.5-BLO15-EMP1.5}. \\
  Desfase termico & $5,2$ h & Retraso aproximado entre el maximo de temperatura exterior y el maximo en la cara interior del componente. \\
  Amortiguamiento de amplitud & $1,8$ & Relacion entre la oscilacion exterior y la oscilacion interior; un valor mayor indica mejor reduccion de la onda termica. \\
  TAD & $0,544$ & Indicador adimensional de reduccion dinamica de la amplitud termica. \\
  Capacidad de almacenamiento del componente & $289\,\mathrm{kJ/m^2K}$ & Capacidad termica total asociada al conjunto de capas del muro. \\
  Capacidad termica de capas interiores & $77\,\mathrm{kJ/m^2K}$ & Porcion de capacidad termica mas cercana al ambiente interior. \\
\end{longtable}
}
\fuente{Elaboracion propia a partir de los resultados de comportamiento dinamico del componente de muro exterior.}
